% =====================================================================
%  Thème 1 — Conversions & cercle trigonométrique — NIVEAU FACILE
%  Exercices
% =====================================================================
\documentclass[11pt,a4paper]{article}
\usepackage[utf8]{inputenc}
\usepackage[T1]{fontenc}
\usepackage{lmodern}
\usepackage[french]{babel}
\usepackage{amsmath,amssymb,mathtools}
\usepackage{icomma}
\usepackage{xcolor}
\usepackage{colortbl}
\usepackage{tikz}
\usepackage[most]{tcolorbox}
\usepackage{enumitem}
\usepackage{array}
\usepackage[a4paper,margin=2.2cm]{geometry}
\usetikzlibrary{arrows.meta,calc}

\definecolor{primary}{RGB}{214,51,108}
\definecolor{primarydark}{RGB}{140,20,64}

\DeclareMathOperator{\tg}{tg}
\DeclareMathOperator{\cotg}{cotg}
\newcommand{\degr}{^{\circ}}
\newcommand{\Z}{\mathbb{Z}}

\newtcolorbox{consigne}{enhanced,breakable,colback=primary!5,colframe=primary,
  boxrule=0.8pt,arc=2pt,left=3mm,right=3mm,top=2mm,bottom=2mm}

\newcounter{exo}
\newcommand{\exo}[1]{\medskip\refstepcounter{exo}%
  \noindent{\large\bfseries\color{primarydark}Exercice \theexo}%
  \ \textcolor{primary}{\textbullet}\ \textbf{#1}\par\nopagebreak\smallskip}

\setlength{\parindent}{0pt}
\pagestyle{empty}

\begin{document}

\begin{center}
{\LARGE\bfseries\color{primary}Thème 1 — Conversions \& cercle trigonométrique}\\[2pt]
{\color{primarydark}\textsc{Niveau Facile — Exercices}}
\end{center}
\hrule height 1pt
\medskip

\begin{consigne}
\textbf{Objectif :} rendre automatiques les conversions, la lecture des quadrants, des signes et des
valeurs remarquables. \textbf{Tout se fait sans calculatrice.} Donne des valeurs \emph{exactes}
(fractions, $\sqrt{\ }$, $\pi$).
\end{consigne}

\exo{Des degrés vers les radians}
Convertis chaque angle en radians, sous forme d'une fraction de $\pi$ simplifiée.
\[
\textbf{(a)}\ 45\degr \qquad
\textbf{(b)}\ 90\degr \qquad
\textbf{(c)}\ 150\degr \qquad
\textbf{(d)}\ 210\degr \qquad
\textbf{(e)}\ 300\degr
\]

\exo{Des radians vers les degrés}
Convertis chaque angle en degrés.
\[
\textbf{(a)}\ \frac{\pi}{6} \qquad
\textbf{(b)}\ \frac{3\pi}{4} \qquad
\textbf{(c)}\ \frac{5\pi}{3} \qquad
\textbf{(d)}\ \frac{7\pi}{6} \qquad
\textbf{(e)}\ 2\pi
\]

\exo{Dans quel quadrant ?}
Indique le quadrant (I, II, III ou IV) auquel appartient chaque angle.
\[
\textbf{(a)}\ 200\degr \qquad
\textbf{(b)}\ 85\degr \qquad
\textbf{(c)}\ \frac{5\pi}{4} \qquad
\textbf{(d)}\ \frac{7\pi}{4} \qquad
\textbf{(e)}\ \frac{2\pi}{3}
\]

\exo{Quel signe ? (sans calculer la valeur)}
Pour chaque nombre trigonométrique, donne uniquement son \textbf{signe} ($+$ ou $-$) en te servant du
quadrant.
\[
\textbf{(a)}\ \cos 160\degr \qquad
\textbf{(b)}\ \sin 250\degr \qquad
\textbf{(c)}\ \tg 100\degr \qquad
\textbf{(d)}\ \cos 320\degr \qquad
\textbf{(e)}\ \sin 40\degr
\]

\exo{Valeurs remarquables du premier quadrant}
Donne la valeur exacte, en t'aidant du tableau des valeurs remarquables.
\[
\textbf{(a)}\ \cos\frac{\pi}{3} \qquad
\textbf{(b)}\ \sin\frac{\pi}{4} \qquad
\textbf{(c)}\ \tg\frac{\pi}{3} \qquad
\textbf{(d)}\ \sin\frac{\pi}{6} \qquad
\textbf{(e)}\ \cos\frac{\pi}{2}
\]

\end{document}
